% ============================================================================
% CSCE 763: Digital Image Processing - Homework 5 Solutions
% ============================================================================

\documentclass[12pt]{article}

% ============================================================================
% PACKAGE IMPORTS
% ============================================================================
\usepackage{amsmath}
\usepackage{graphicx}
\usepackage{amssymb}
\usepackage{tikz}
\usepackage[margin=1in]{geometry}
\usepackage{setspace}
\usepackage{xcolor}
\usepackage{enumitem}
\usepackage{tcolorbox}
\usepackage{fancyhdr}
\usepackage{booktabs}
\usepackage{array}
\usepackage{colortbl}
\usepackage{float}

% ============================================================================
% HEADER AND FOOTER SETUP
% ============================================================================
\pagestyle{fancy}
\setlength{\headheight}{14.5pt}
\fancyhf{}

\fancyhead[L]{Page \thepage}
\fancyhead[C]{CSCE 763}
\fancyhead[R]{J.C. Vaught}

\renewcommand{\headrulewidth}{0pt}

% ============================================================================
% TIKZ LIBRARIES
% ============================================================================
\usetikzlibrary{shapes.geometric, arrows.meta, positioning, patterns, calc}

% --- USC Brand Colors ---
\definecolor{USC_Garnet}{HTML}{73000A}
\definecolor{USC_Sandstorm}{HTML}{FFF2E3}
\definecolor{USC_Black90}{HTML}{363636}
\definecolor{USC_Black70}{HTML}{5C5C5C}
\definecolor{USC_Black50}{HTML}{A2A2A2}
\definecolor{USC_Black10}{HTML}{ECECEC}
\definecolor{USC_Honeycomb}{HTML}{A49137}
\definecolor{USC_Atlantic}{HTML}{466A9F}
\definecolor{USC_Congaree}{HTML}{1F414D}
\definecolor{USC_Rose}{HTML}{CC2E40}
\definecolor{USC_Grass}{HTML}{CED318}
\definecolor{USC_Horseshoe}{HTML}{65780B}
\definecolor{outputbg}{RGB}{245,245,245}

% ============================================================================
% CUSTOM COLORED BOXES
% ============================================================================
\newtcolorbox{stepbox}[1][Step]{
  colback=white,
  colframe=USC_Garnet,
  fonttitle=\bfseries,
  title={#1},
  sharp corners,
  colbacktitle=USC_Garnet,
  coltitle=white,
}

\newtcolorbox{resultsbox}{
  colback=white,
  colframe=USC_Honeycomb,
  fonttitle=\bfseries,
  title=Results,
  sharp corners,
  colbacktitle=USC_Honeycomb,
  coltitle=white,
}

% ============================================================================
% CUSTOM QUESTION AND PART COMMANDS
% ============================================================================
\newcommand{\question}[1]{%
  \clearpage
  \vspace{0.5cm}
  {\noindent\normalsize \textbf{#1}}
  \vspace{0.2cm}
  \hrule
  \vspace{0.1cm}
  \hrule
  \vspace{0.3cm}
}

\newcommand{\questionpart}[1]{%
  \clearpage
  \vspace{0.5cm}
  {\noindent\normalsize \textbf{#1}}
  \vspace{0.2cm}
  \hrule
  \vspace{0.1cm}
  \hrule
  \vspace{0.3cm}
}

% ============================================================================
% DOCUMENT INFORMATION
% ============================================================================

\title{CSCE 763: Digital Image Processing \\ Homework \#5}
\author{Instructor: Dr. Yan Tong \\ University of South Carolina \\ \textit{Solutions by: J.C. Vaught}}
\date{Due: 1:15 pm EST, Thursday, April 9}

% ============================================================================
% BEGIN DOCUMENT
% ============================================================================

\begin{document}

\maketitle
\setlength{\parindent}{0pt}

\setlist[enumerate,1]{label=\arabic*.}
\setlist[enumerate,2]{label=(\alph*)}

% ============================================================================
% TABLE OF CONTENTS
% ============================================================================

\begin{center}
\begin{tcolorbox}[colback=white, colframe=USC_Black90, colbacktitle=USC_Black10,
                  coltitle=USC_Black90, sharp corners, boxrule=1pt,
                  title=\Large\bfseries Table of Contents]
\begin{tabular}{p{0.82\textwidth}r}
\textbf{Problem 1: Image Entropy and Compression (25 pts)} \dotfill & \textbf{\pageref{quest:1}} \\[0.3em]
\textbf{Problem 2: Arithmetic Decoding (25 pts)} \dotfill & \textbf{\pageref{quest:2}} \\[0.3em]
\textbf{Problem 3: Line Break Detection Masks (25 pts)} \dotfill & \textbf{\pageref{quest:3}} \\[0.3em]
\textbf{Problem 4: Sobel Mask Decomposition (25 pts)} \dotfill & \textbf{\pageref{quest:4}} \\
\end{tabular}
\vspace{0.2cm}
\end{tcolorbox}
\end{center}

\vspace{0.5cm}

% ============================================================================
% PROBLEM 1: Image Entropy and Compression
% ============================================================================

\question{Problem 1: Image Entropy and Compression (25 pts)}\label{quest:1}

Consider a simple $4 \times 16$, 8-bit image:

\vspace{0.3cm}

\begin{center}
\ttfamily
\begin{tabular}{cccccccccccccccc}
22 & 22 & 21 & 94 & 93 & 94 & 166 & 167 & 166 & 167 & 253 & 254 & 254 & 255 & 253 & 55 \\
20 & 22 & 21 & 22 & 94 & 95 & 93 & 166 & 167 & 166 & 254 & 253 & 255 & 255 & 254 & 54 \\
21 & 21 & 20 & 22 & 95 & 95 & 95 & 165 & 166 & 166 & 255 & 254 & 254 & 254 & 253 & 54 \\
21 & 21 & 22 & 95 & 95 & 94 & 166 & 167 & 165 & 166 & 254 & 255 & 253 & 253 & 255 & 55 \\
\end{tabular}
\end{center}
\rmfamily

\vspace{0.3cm}

\begin{enumerate}[label=(\alph*)]
\item Compute the entropy of the image. (5 pts)
\item Compress the image using Huffman coding and report the compression ratio using Huffman coding. (10 pts)
\item Use another data compression algorithm shown in the lecture to obtain a higher compression ratio. (10 pts)
\end{enumerate}

% --- Part (a): Entropy ---

\vspace{0.4cm}

\begin{stepbox}[Part (a): Compute Entropy]
The image has $4 \times 16 = 64$ pixels. Counting the frequency of each distinct intensity value yields 14 unique symbols.

\vspace{0.3cm}

\begin{center}
\renewcommand{\arraystretch}{1.1}
\begin{tabular}{|c|c|c||c|c|c|}
\hline
\textbf{Value} & \textbf{Count} & $\boldsymbol{p_i}$ & \textbf{Value} & \textbf{Count} & $\boldsymbol{p_i}$ \\
\hline
20  & 2 & 2/64  & 95  & 6 & 6/64 \\
21  & 6 & 6/64  & 165 & 2 & 2/64 \\
22  & 6 & 6/64  & 166 & 8 & 8/64 \\
54  & 2 & 2/64  & 167 & 4 & 4/64 \\
55  & 2 & 2/64  & 253 & 6 & 6/64 \\
93  & 2 & 2/64  & 254 & 8 & 8/64 \\
94  & 4 & 4/64  & 255 & 6 & 6/64 \\
\hline
\end{tabular}
\end{center}

\vspace{0.3cm}

Group the symbols by probability to compute entropy $H = -\sum p_i \log_2 p_i$:

\begin{itemize}
\item 5 symbols with $p=2/64$: each contributes $-\frac{2}{64}\log_2\frac{2}{64} = \frac{5}{32}$. Subtotal: $5 \times \frac{5}{32} = \frac{25}{32} = 0.78125$.
\item 5 symbols with $p=6/64$: each contributes $-\frac{6}{64}\log_2\frac{6}{64} = \frac{3}{32}(5-\log_2 3)$. Subtotal: $5 \times 0.32016 = 1.60078$.
\item 2 symbols with $p=4/64$: each contributes $-\frac{4}{64}\log_2\frac{4}{64} = \frac{4}{16} = 0.25$. Subtotal: $2 \times 0.25 = 0.50$.
\item 2 symbols with $p=8/64$: each contributes $-\frac{8}{64}\log_2\frac{8}{64} = \frac{3}{8} = 0.375$. Subtotal: $2 \times 0.375 = 0.75$.
\end{itemize}
\end{stepbox}

\begin{resultsbox}
\[
\boxed{H = 0.78125 + 1.60078 + 0.50 + 0.75 = 3.632 \text{ bits/pixel}}
\]
\end{resultsbox}

% --- Part (b): Huffman Coding ---

\vspace{0.4cm}

\begin{stepbox}[Part (b), Step 1: Build Huffman Tree]
Build a Huffman tree by repeatedly combining the two lowest-frequency symbols.

Sort symbols by frequency (ascending): 20(2), 54(2), 55(2), 93(2), 165(2), 94(4), 167(4), 21(6), 22(6), 95(6), 253(6), 255(6), 166(8), 254(8).

\begin{enumerate}[label=\arabic*.]
\item Combine 20(2) + 54(2) $\to$ $N_1$(4)
\item Combine 55(2) + 93(2) $\to$ $N_2$(4)
\item Combine 165(2) + $N_1$(4) $\to$ $N_3$(6)
\item Combine $N_2$(4) + 94(4) $\to$ $N_4$(8)
\item Combine 167(4) + $N_3$(6) $\to$ $N_5$(10)
\item Combine 21(6) + 22(6) $\to$ $N_6$(12)
\item Combine 95(6) + 253(6) $\to$ $N_7$(12)
\item Combine 255(6) + $N_4$(8) $\to$ $N_8$(14)
\item Combine 166(8) + 254(8) $\to$ $N_9$(16)
\item Combine $N_5$(10) + $N_6$(12) $\to$ $N_{10}$(22)
\item Combine $N_7$(12) + $N_8$(14) $\to$ $N_{11}$(26)
\item Combine $N_9$(16) + $N_{10}$(22) $\to$ $N_{12}$(38)
\item Combine $N_{11}$(26) + $N_{12}$(38) $\to$ Root(64)
\end{enumerate}
\end{stepbox}

\begin{stepbox}[Part (b), Step 2: Assign Huffman Codes]
Traversing the tree from root to each leaf, assigning 0 for left branches and 1 for right branches, yields the following code table. Because several symbols share the same frequency, other optimal Huffman trees are also possible. The exact bit patterns can vary, but the total coded length remains the same.

\begin{center}
\renewcommand{\arraystretch}{1.15}
\begin{tabular}{|c|c|c|c|c|}
\hline
\textbf{Symbol} & \textbf{Count} & \textbf{Code} & \textbf{Bits} & \textbf{Total Bits} \\
\hline
166  & 8 & 100    & 3 & 24 \\
254  & 8 & 101    & 3 & 24 \\
95   & 6 & 000    & 3 & 18 \\
253  & 6 & 001    & 3 & 18 \\
255  & 6 & 010    & 3 & 18 \\
21   & 6 & 1110   & 4 & 24 \\
22   & 6 & 1111   & 4 & 24 \\
94   & 4 & 0111   & 4 & 16 \\
167  & 4 & 1100   & 4 & 16 \\
55   & 2 & 01100  & 5 & 10 \\
93   & 2 & 01101  & 5 & 10 \\
165  & 2 & 11010  & 5 & 10 \\
20   & 2 & 110110 & 6 & 12 \\
54   & 2 & 110111 & 6 & 12 \\
\hline
\multicolumn{4}{|r|}{\textbf{Grand Total:}} & \textbf{236} \\
\hline
\end{tabular}
\end{center}
\end{stepbox}

\begin{resultsbox}
Average code length: $236/64 = 3.6875$ bits/pixel (close to entropy of 3.632).

Original image: $64 \times 8 = 512$ bits. Huffman-coded: 236 bits.
\[
\boxed{\text{Compression Ratio} = \frac{512}{236} \approx 2.169}
\]
\end{resultsbox}

% --- Part (c): Higher Compression ---

\vspace{0.4cm}

\begin{stepbox}[Part (c), Step 1: Predictive Coding Approach]
Huffman coding on the raw pixel values ignores spatial correlation between neighboring pixels. A higher compression ratio can be achieved by first applying \textit{differential pulse code modulation} (DPCM) to exploit the strong inter-pixel redundancy, then Huffman coding the prediction errors.

\textbf{Prediction scheme:} The first pixel is sent directly (8 bits). Each subsequent pixel in a row is predicted by its left neighbor. The first pixel of rows 2--4 is predicted by the pixel directly above it. The transmitted quantity is the prediction error $e_n = x_n - \hat{x}_n$.
\end{stepbox}

\begin{stepbox}[Part (c), Step 2: Compute Prediction Errors]
Applying the predictor to each row yields 63 error values.

{\small
\textbf{Row 1} (cols 2--16): $0, -1, 73, -1, 1, 72, 1, -1, 1, 86, 1, 0, 1, -2, -198$

\textbf{Row 2} (col 1 from above, then cols 2--16):

$-2, 2, -1, 1, 72, 1, -2, 73, 1, -1, 88, -1, 2, 0, -1, -200$

\textbf{Row 3}: $1, 0, -1, 2, 73, 0, 0, 70, 1, 0, 89, -1, 0, 0, -1, -199$

\textbf{Row 4}: $0, 0, 1, 73, 0, -1, 72, 1, -2, 1, 88, 1, -2, 0, 2, -200$
}

The errors cluster heavily around 0, with large jumps only at intensity-group boundaries.
\end{stepbox}

\begin{stepbox}[Part (c), Step 3: Huffman Code the Errors]
Building a Huffman tree on the 14 distinct error values yields the following code table.

\begin{center}
\renewcommand{\arraystretch}{1.15}
\begin{tabular}{|c|c|c|c|c|}
\hline
\textbf{Error} & \textbf{Count} & \textbf{Code} & \textbf{Bits} & \textbf{Total Bits} \\
\hline
1    & 14 & 01      & 2 & 28 \\
0    & 13 & 00      & 2 & 26 \\
$-1$ & 11 & 111     & 3 & 33 \\
$-2$ & 5  & 1101    & 4 & 20 \\
2    & 4  & 1010    & 4 & 16 \\
73   & 4  & 1011    & 4 & 16 \\
72   & 3  & 1000    & 4 & 12 \\
88   & 2  & 10011   & 5 & 10 \\
$-200$& 2 & 11000   & 5 & 10 \\
89   & 1  & 100100  & 6 & 6  \\
$-198$& 1 & 100101  & 6 & 6  \\
$-199$& 1 & 110010  & 6 & 6  \\
70   & 1  & 1100110 & 7 & 7  \\
86   & 1  & 1100111 & 7 & 7  \\
\hline
\multicolumn{4}{|r|}{\textbf{Error bits:}} & \textbf{203} \\
\hline
\end{tabular}
\end{center}

Total bits $= 8 \text{ (first pixel)} + 203 \text{ (Huffman-coded errors)} = 211$ bits.
\end{stepbox}

\begin{resultsbox}
The prediction errors have entropy $H_e \approx 3.186$ bits/error symbol, lower than the raw-pixel entropy of $3.632$ bits/pixel. Including the first uncoded pixel, the overall average rate is
\[
\frac{211}{64} \approx 3.297 \text{ bits/pixel}.
\]
This confirms that DPCM removes significant inter-pixel redundancy before Huffman coding.

\[
\boxed{\text{Compression Ratio} = \frac{512}{211} \approx 2.427}
\]

This exceeds the Huffman-only ratio of 2.169 because predictive coding exploits the spatial correlation that symbol-by-symbol Huffman coding cannot.
\end{resultsbox}


% ============================================================================
% PROBLEM 2: Arithmetic Decoding
% ============================================================================

\question{Problem 2: Arithmetic Decoding (25 pts)}\label{quest:2}

The arithmetic decoding process is the reverse of the encoding procedure. Decode the message
$0.222355$ given the coding model. Assume that ``!'' is the symbol indicating the end of message.
Show the decoding process.

\vspace{0.3cm}

\begin{center}
\begin{tabular}{cc}
\toprule
\textbf{Symbol} & \textbf{Probability} \\
\midrule
A & 0.25 \\
B & 0.15 \\
C & 0.1 \\
D & 0.3 \\
E & 0.1 \\
! & 0.1 \\
\bottomrule
\end{tabular}
\end{center}

\vspace{0.4cm}

\begin{stepbox}[Step 1: Establish Cumulative Probability Ranges]
Assign cumulative subintervals to each symbol in the order given.

\begin{center}
\renewcommand{\arraystretch}{1.15}
\begin{tabular}{|c|c|c|}
\hline
\textbf{Symbol} & \textbf{Probability} & \textbf{Range} \\
\hline
A & 0.25 & $[0,\; 0.25)$ \\
B & 0.15 & $[0.25,\; 0.40)$ \\
C & 0.10 & $[0.40,\; 0.50)$ \\
D & 0.30 & $[0.50,\; 0.80)$ \\
E & 0.10 & $[0.80,\; 0.90)$ \\
! & 0.10 & $[0.90,\; 1.00)$ \\
\hline
\end{tabular}
\end{center}
\end{stepbox}

\begin{stepbox}[Step 2: Iterative Decoding]
Start with interval $[\text{low}, \text{high}) = [0, 1)$ and code value $v = 0.222355$.

\vspace{0.3cm}

\textbf{Iteration 1:} $[\text{low}, \text{high}) = [0,\; 1)$, width $= 1$

Subintervals:
\begin{center}
{\small A: $[0,\;0.25)$,\quad B: $[0.25,\;0.40)$,\quad C: $[0.40,\;0.50)$,\quad D: $[0.50,\;0.80)$,\quad E: $[0.80,\;0.90)$,\quad !: $[0.90,\;1.00)$}
\end{center}
$v = 0.222355$ falls in $[0,\;0.25)$ $\Rightarrow$ \textbf{Output: A}

Update: $\text{low} = 0$, $\text{high} = 0.25$

\vspace{0.3cm}

\textbf{Iteration 2:} $[\text{low}, \text{high}) = [0,\; 0.25)$, width $= 0.25$

Subintervals:
\begin{center}
{\small A: $[0,\;0.0625)$,\; B: $[0.0625,\;0.1)$,\; C: $[0.1,\;0.125)$,\; D: $[0.125,\;0.2)$,\; E: $[0.2,\;0.225)$,\; !: $[0.225,\;0.25)$}
\end{center}
$v = 0.222355$ falls in $[0.2,\;0.225)$ $\Rightarrow$ \textbf{Output: E}

Update: $\text{low} = 0.2$, $\text{high} = 0.225$

\vspace{0.3cm}

\textbf{Iteration 3:} $[\text{low}, \text{high}) = [0.2,\; 0.225)$, width $= 0.025$

Subintervals:
\begin{center}
{\small A: $[0.2,\;0.20625)$,\; B: $[0.20625,\;0.21)$,\; C: $[0.21,\;0.2125)$,\; D: $[0.2125,\;0.22)$}

{\small E: $[0.22,\;0.2225)$,\; !: $[0.2225,\;0.225)$}
\end{center}
$v = 0.222355$ falls in $[0.22,\;0.2225)$ $\Rightarrow$ \textbf{Output: E}

Update: $\text{low} = 0.22$, $\text{high} = 0.2225$

\vspace{0.3cm}

\textbf{Iteration 4:} $[\text{low}, \text{high}) = [0.22,\; 0.2225)$, width $= 0.0025$

Subintervals:
\begin{center}
{\small A: $[0.22,\;0.220625)$,\; B: $[0.220625,\;0.221)$,\; C: $[0.221,\;0.22125)$}

{\small D: $[0.22125,\;0.222)$,\; E: $[0.222,\;0.22225)$,\; !: $[0.22225,\;0.2225)$}
\end{center}
$v = 0.222355$ falls in $[0.22225,\;0.2225)$ $\Rightarrow$ \textbf{Output: !}

End-of-message symbol received. Stop decoding.
\end{stepbox}

\begin{stepbox}[Verification: Encode ``AEE!'' to Confirm]
\begin{center}
\renewcommand{\arraystretch}{1.3}
\begin{tabular}{|c|c|c|c|}
\hline
\textbf{Step} & \textbf{Symbol} & \textbf{New Low} & \textbf{New High} \\
\hline
Init &  & 0 & 1 \\
1 & A & $0 + 1 \times 0 = 0$ & $0 + 1 \times 0.25 = 0.25$ \\
2 & E & $0 + 0.25 \times 0.8 = 0.2$ & $0 + 0.25 \times 0.9 = 0.225$ \\
3 & E & $0.2 + 0.025 \times 0.8 = 0.22$ & $0.2 + 0.025 \times 0.9 = 0.2225$ \\
4 & ! & $0.22 + 0.0025 \times 0.9 = 0.22225$ & $0.22 + 0.0025 \times 1.0 = 0.2225$ \\
\hline
\end{tabular}
\end{center}
Final interval: $[0.22225,\; 0.2225)$. The code value $0.222355 \in [0.22225,\; 0.2225)$. $\checkmark$
\end{stepbox}

\begin{resultsbox}
\[
\boxed{\text{Decoded message: } \mathbf{A\;E\;E\;!}}
\]
\end{resultsbox}


% ============================================================================
% PROBLEM 3: Line Break Detection Masks
% ============================================================================

\question{Problem 3: Line Break Detection Masks (25 pts)}\label{quest:3}

A binary image contains straight lines oriented horizontally, vertically, at $45^\circ$, and at $-45^\circ$.
Develop a set of $3 \times 3$ masks that can be used to detect 1-pixel breaks in these lines. For example,
you may want to detect a pattern like $1\ 1\ 1\ 0\ 1\ 1\ 1$ in the vertical direction. Assume that the
intensities of the lines and background are represented by 1 and 0, respectively.

\vspace{0.4cm}

\begin{stepbox}[Step 1: Break Pattern Analysis]
A 1-pixel break in a line means the center pixel is 0 (the gap) while the two nearest pixels along the line direction are 1 (the line continues on both sides). A good detection mask should produce a strong positive response at break locations and zero response at normal line points and background.

Consider the $3 \times 3$ neighborhood centered on the break pixel. The mask assigns positive weights to the positions where line pixels are expected (along the line direction, adjacent to the gap) and a negative weight at the center (which should be 0 for a break).
\end{stepbox}

\begin{stepbox}[Step 2: Design the Four Masks]

\textbf{Horizontal line break} --- the line runs left-to-right, break at center:
\[
\text{Pattern: }
\begin{bmatrix} 0 & 0 & 0 \\ 1 & 0 & 1 \\ 0 & 0 & 0 \end{bmatrix}
\qquad
\mathbf{M}_H =
\begin{bmatrix} 0 & 0 & 0 \\ 1 & -2 & 1 \\ 0 & 0 & 0 \end{bmatrix}
\]

\textbf{Vertical line break} --- the line runs top-to-bottom, break at center:
\[
\text{Pattern: }
\begin{bmatrix} 0 & 1 & 0 \\ 0 & 0 & 0 \\ 0 & 1 & 0 \end{bmatrix}
\qquad
\mathbf{M}_V =
\begin{bmatrix} 0 & 1 & 0 \\ 0 & -2 & 0 \\ 0 & 1 & 0 \end{bmatrix}
\]

\textbf{$45^\circ$ line break} --- the line runs from top-right to bottom-left, break at center:
\[
\text{Pattern: }
\begin{bmatrix} 0 & 0 & 1 \\ 0 & 0 & 0 \\ 1 & 0 & 0 \end{bmatrix}
\qquad
\mathbf{M}_{45} =
\begin{bmatrix} 0 & 0 & 1 \\ 0 & -2 & 0 \\ 1 & 0 & 0 \end{bmatrix}
\]

\textbf{$-45^\circ$ line break} --- the line runs from top-left to bottom-right, break at center:
\[
\text{Pattern: }
\begin{bmatrix} 1 & 0 & 0 \\ 0 & 0 & 0 \\ 0 & 0 & 1 \end{bmatrix}
\qquad
\mathbf{M}_{-45} =
\begin{bmatrix} 1 & 0 & 0 \\ 0 & -2 & 0 \\ 0 & 0 & 1 \end{bmatrix}
\]
\end{stepbox}

\begin{stepbox}[Step 3: Response Verification]
Each mask gives the following responses when centered on different pixel types.

\begin{center}
\renewcommand{\arraystretch}{1.25}
\small
\begin{tabular}{|l|c|c|c|}
\hline
\textbf{Location} & \textbf{Center} & \textbf{Line Neighbors} & \textbf{Response} \\
\hline
Break point  & 0 & both = 1 & $1+(-2)(0)+1 = \mathbf{2}$ \\
Normal line  & 1 & both = 1 & $1+(-2)(1)+1 = \mathbf{0}$ \\
Background   & 0 & both = 0 & $0+(-2)(0)+0 = \mathbf{0}$ \\
Endpoint     & 1 & 1 and 0  & $1+(-2)(1)+0 = \mathbf{-1}$ \\
\hline
\end{tabular}
\normalsize
\end{center}

A threshold of $>0$ (or specifically $=2$) on the mask response cleanly isolates 1-pixel break locations. The masks produce the maximum response of 2 exclusively at break points, zero at normal line points and background, and $-1$ at endpoints.
\end{stepbox}

\begin{resultsbox}
The four $3 \times 3$ break-detection masks are:
\[
\boxed{
\mathbf{M}_H = \begin{bmatrix} 0 & 0 & 0 \\ 1 & -2 & 1 \\ 0 & 0 & 0 \end{bmatrix}, \quad
\mathbf{M}_V = \begin{bmatrix} 0 & 1 & 0 \\ 0 & -2 & 0 \\ 0 & 1 & 0 \end{bmatrix}
}
\]
\[
\boxed{
\mathbf{M}_{45} = \begin{bmatrix} 0 & 0 & 1 \\ 0 & -2 & 0 \\ 1 & 0 & 0 \end{bmatrix}, \quad
\mathbf{M}_{-45} = \begin{bmatrix} 1 & 0 & 0 \\ 0 & -2 & 0 \\ 0 & 0 & 1 \end{bmatrix}
}
\]
A response of 2 at any pixel indicates a 1-pixel break in the corresponding line orientation.
\end{resultsbox}


% ============================================================================
% PROBLEM 4: Sobel Mask Decomposition
% ============================================================================

\question{Problem 4: Sobel Mask Decomposition (25 pts)}\label{quest:4}

The results obtained by a single pass through an image of some 2D masks can be achieved also
by two passes using 1D masks for improving efficiency. Show that the response of Sobel masks

\begin{align}
\mathbf{S}_x &=
\begin{bmatrix}
-1 & -2 & -1 \\
0 & 0 & 0 \\
1 & 2 & 1
\end{bmatrix} \\[1em]
\mathbf{S}_y &=
\begin{bmatrix}
-1 & 0 & 1 \\
-2 & 0 & 2 \\
-1 & 0 & 1
\end{bmatrix}
\end{align}

can be implemented similarly by one pass of a vertical differencing mask

\begin{align}
\mathbf{d}_v &=
\begin{bmatrix}
-1 \\
0 \\
1
\end{bmatrix}
\end{align}

followed by a horizontal smoothing mask

\begin{align}
\mathbf{s}_h &=
\begin{bmatrix} 1 & 2 & 1 \end{bmatrix}
\end{align}

\vspace{0.4cm}

\begin{stepbox}[Step 1: Separability of 2D Convolution]
A $3 \times 3$ kernel $\mathbf{H}$ is \textit{separable} if it can be expressed as the outer product of a column vector and a row vector:
\[
\mathbf{H} = \mathbf{c} \cdot \mathbf{r}^T
\]
When separable, the 2D convolution $f \ast \mathbf{H}$ can be decomposed into two sequential 1D convolutions:
\[
f \ast \mathbf{H} = (f \ast_{\text{cols}} \mathbf{c}) \ast_{\text{rows}} \mathbf{r}
\]
This reduces the computational cost from $O(m^2)$ multiplications per pixel (for an $m \times m$ kernel) to $O(2m)$.
\end{stepbox}

\begin{stepbox}[Step 2: Decompose $\mathbf{S}_x$]
Compute the outer product of $\mathbf{d}_v$ and $\mathbf{s}_h$:
\[
\mathbf{d}_v \cdot \mathbf{s}_h
=
\begin{bmatrix} -1 \\ 0 \\ 1 \end{bmatrix}
\begin{bmatrix} 1 & 2 & 1 \end{bmatrix}
=
\begin{bmatrix}
(-1)(1) & (-1)(2) & (-1)(1) \\
(0)(1)  & (0)(2)  & (0)(1)  \\
(1)(1)  & (1)(2)  & (1)(1)
\end{bmatrix}
\]
\[
=
\begin{bmatrix}
-1 & -2 & -1 \\
0  & 0  & 0  \\
1  & 2  & 1
\end{bmatrix}
= \mathbf{S}_x \quad \checkmark
\]
Therefore, the 2D convolution with $\mathbf{S}_x$ can be computed by first convolving each column with $\mathbf{d}_v$ (vertical differencing), then convolving each row of the result with $\mathbf{s}_h$ (horizontal smoothing).
\end{stepbox}

\begin{stepbox}[Step 3: Decompose $\mathbf{S}_y$]
Define the complementary 1D masks: vertical smoothing $\mathbf{s}_v = \begin{bmatrix} 1 \\ 2 \\ 1 \end{bmatrix}$ and horizontal differencing $\mathbf{d}_h = \begin{bmatrix} -1 & 0 & 1 \end{bmatrix}$.

Compute the outer product:
\[
\mathbf{s}_v \cdot \mathbf{d}_h
=
\begin{bmatrix} 1 \\ 2 \\ 1 \end{bmatrix}
\begin{bmatrix} -1 & 0 & 1 \end{bmatrix}
=
\begin{bmatrix}
(1)(-1) & (1)(0) & (1)(1) \\
(2)(-1) & (2)(0) & (2)(1) \\
(1)(-1) & (1)(0) & (1)(1)
\end{bmatrix}
\]
\[
=
\begin{bmatrix}
-1 & 0 & 1 \\
-2 & 0 & 2 \\
-1 & 0 & 1
\end{bmatrix}
= \mathbf{S}_y \quad \checkmark
\]
Therefore, $\mathbf{S}_y$ is computed by first convolving each column with $\mathbf{s}_v$ (vertical smoothing), then convolving each row with $\mathbf{d}_h$ (horizontal differencing).
\end{stepbox}

\begin{stepbox}[Step 4: Equivalence of the Two-Pass Implementation]
Both Sobel masks are separable products of the same two primitive 1D operations --- a 3-point central difference $[-1,\; 0,\; 1]$ and a 3-point binomial smoothing $[1,\; 2,\; 1]$ --- applied in perpendicular directions.

\begin{center}
\renewcommand{\arraystretch}{1.3}
\begin{tabular}{|c|c|c|}
\hline
\textbf{Sobel Mask} & \textbf{Pass 1 (columns)} & \textbf{Pass 2 (rows)} \\
\hline
$\mathbf{S}_x$ (vertical edges) & $\mathbf{d}_v = [-1;\; 0;\; 1]$ (difference) & $\mathbf{s}_h = [1,\; 2,\; 1]$ (smooth) \\
\hline
$\mathbf{S}_y$ (horizontal edges) & $\mathbf{s}_v = [1;\; 2;\; 1]$ (smooth) & $\mathbf{d}_h = [-1,\; 0,\; 1]$ (difference) \\
\hline
\end{tabular}
\end{center}

Since convolution is commutative and associative, the order of the two 1D passes can also be swapped.
\end{stepbox}

\begin{resultsbox}
Both Sobel masks are separable. Their responses are equivalent to two sequential 1D operations.

\[
\boxed{\mathbf{S}_x = \mathbf{d}_v \cdot \mathbf{s}_h =
\begin{bmatrix} -1 \\ 0 \\ 1 \end{bmatrix}
\begin{bmatrix} 1 & 2 & 1 \end{bmatrix}}
\qquad
\boxed{\mathbf{S}_y = \mathbf{s}_v \cdot \mathbf{d}_h =
\begin{bmatrix} 1 \\ 2 \\ 1 \end{bmatrix}
\begin{bmatrix} -1 & 0 & 1 \end{bmatrix}}
\]

This reduces the cost per pixel from 9 multiplications to 6 (two passes of 3 each).
\end{resultsbox}


% ============================================================================
% END OF DOCUMENT
% ============================================================================

\end{document}
